\documentclass{standalone}
\usepackage{ctex}
\usepackage{tikz}
\usetikzlibrary{
    shapes.geometric,
    arrows.meta,
    positioning,
    calc
}

\tikzset{
    base/.style = {
        draw,
        align=center,
        minimum width=3.2cm,
        minimum height=1cm,
    },
    terminal/.style = {
        base,
        rectangle,
        rounded corners=0.5cm,
        minimum width=2cm,
        minimum height=1cm,
    },
    io/.style = {
        base,
        trapezium,
        trapezium left angle=60,
        trapezium right angle=120,
        trapezium stretches=true,
    },
    process/.style = {
        base,
        rectangle,
    },
    decision/.style = {
        base,
        diamond,
        aspect=2,
        inner sep=1pt,
    },
    arrow/.style = {
        ->,
        thick,
        >={Stealth},
    },
    line/.style = {
        thick,
    },
}

\begin{document}
\begin{tikzpicture}[node distance=1.4cm]

    % 统一控制竖向边到边距离
    \def\vGap{0.8cm}

    \node [terminal] (start) {开始};
    \node [decision] (role) [below=of start] {用户类型};

    \draw [arrow] (start) -- (role);

    % 用户类型到隐形点：边到点距离为 0.8cm
    \coordinate (rolePoint) at ($(role.south)+(0,-\vGap)$);
    \draw [line] (role.south) -- (rolePoint);

    % 左右分支点
    \coordinate (leftBranch) at ($(rolePoint)+(-3.0cm,0)$);
    \coordinate (rightBranch) at ($(rolePoint)+(3.0cm,0)$);

    % 两个子节点：分支点到节点上边缘距离为 0.8cm
    \node [io] (publicReq) [below=\vGap of leftBranch] {请求已通过友链};
    \node [io] (adminReq) [below=\vGap of rightBranch] {请求友链列表};

    % 隐形点先横着，再竖着连接两个子节点
    \draw [arrow] (rolePoint) -| node [pos=0.25, above] {普通用户} (publicReq.north);
    \draw [arrow] (rolePoint) -| node [pos=0.25, above] {管理员} (adminReq.north);

    % 左侧流程
    \node [process] (show) [below=of publicReq] {展示友链卡片};
    \node [process] (jump) [below=of show] {跳转官方网站};

    \draw [arrow] (publicReq) -- (show);
    \draw [arrow] (show) -- (jump);

    % 右侧流程
    \node [process] (manage) [below=of adminReq] {编辑、审核或删除友链};
    \node [io] (save) [below=of manage] {提交友链保存接口};

    \draw [arrow] (adminReq) -- (manage);
    \draw [arrow] (manage) -- (save);

    % 底部隐形汇合点：节点下边缘到汇合点距离为 0.8cm
    \coordinate (endPoint) at ($(jump.south)!0.5!(save.south)+(0,-\vGap)$);

    % 两个分支先竖着，再横着连接隐形点，不带箭头
    \draw [line] (jump.south) |- (endPoint);
    \draw [line] (save.south) |- (endPoint);

    % 隐形点到结束节点上边缘距离为 0.8cm
    \node [terminal] (stop) [below=\vGap of endPoint] {结束};

    \draw [arrow] (endPoint) -- (stop.north);

\end{tikzpicture}
\end{document}
